\unit{Exposição 236 -- Os esquemas de Picard: propriedades gerais}
\sid{236}
\src{221}
\seminares{14.º ano, 1961/62, n.º 236\\Maio de 1962}
\paperhead{TÉCNICA DE DESCIDA E TEOREMAS DE EXISTÊNCIA EM GEOMETRIA ALGÉBRICA\\
VI. OS ESQUEMAS DE PICARD: PROPRIEDADES GERAIS}%
  {por Alexander GROTHENDIECK}

\fsection{0}{Complementos à exposição anterior ([3], V)\textsuperscript{(1)}}

Houve alguns avanços nas questões de existência dos pré-esquemas de Picard
levantadas em [3], V:

\medskip
\noindent a. (Mumford). Não é verdade, em geral, que, quando
$f:X\longrightarrow S$ é um morfismo projetivo e separável (= plano com fibras
separáveis), exista o pré-esquema de Picard $\operatorname{Pic}_{X/S}$, mesmo
que as fibras de $f$ tenham dimensão 1 e que $S$ seja o espectro de um anel de
valorização discreta completo. Um exemplo é obtido tomando
$S=\operatorname{Spec}\mathbf R[[t]]$ e, para $X$, o subesquema de
$\mathbf P^2_S$ definido pela equação (onde $x,y,z$ são variáveis homogêneas)
\[
x^2+y^2=tz^2,
\]
que representa, portanto, uma cônica que degenera geometricamente em duas
retas concorrentes, embora a fibra especial sobre o corpo $\mathbf R$ seja,
não obstante, irredutível (ela é dada pela equação $x^2+y^2=0$ sobre
$\mathbf R$). Vê-se facilmente que, após a extensão étale
$S'\longrightarrow S$, com $S'=\operatorname{Spec}\mathbf C[[t]]$, existe o
pré-esquema de Picard de $X'/S'$, do qual se obtém uma descrição explícita como
união disjunta de cópias de $\widetilde S$, onde $\widetilde S$ é obtido de
$S$ duplicando infinitas vezes a origem. Verifica-se facilmente que o dado de
descida sobre $\operatorname{Pic}_{X'/S'}$ para $S'\longrightarrow S$ (dado
aqui pela ação do grupo de Galois $\mathbf Z/2\mathbf Z$ de $S'$ sobre $S$)
não é efetivo: o grupo permuta entre si certos pontos duplicados (de modo que
há órbitas que não estão contidas em um aberto afim). Mumford pode demonstrar,
contudo, que, se $f:X\longrightarrow S$ é um morfismo projetivo separável tal
que, para todo $s\in S$, as componentes irredutíveis de $X_s$ são
geometricamente irredutíveis sobre $k(s)$, então
$\operatorname{Pic}_{X/S}$ existe; sua demonstração apoia-se em um refinamento
de seu teorema de passagem ao quociente, cf. [7]. Observe-se, por outro lado,
que continua sendo possível que, sem hipóteses sobre as componentes
irredutíveis das fibras $X_s$, o esquema $\operatorname{Pic}^{\tau}_{X/S}$
(que será introduzido mais adiante) exista mesmo assim.
\src{222}
\noindent b. (Murre). Seja $X$ um esquema próprio sobre um corpo. Então
$\operatorname{Pic}_{X/k}$ existe. A demonstração retoma em parte a
demonstração de Chevalley [1] e apoia-se essencialmente na estrutura de grupo
do funtor de Picard.

Para alguns complementos relativos à teoria dos esquemas de Picard,
especialmente em relação com os esquemas abelianos, será útil consultar [7].
Por fim, uma lacuna notável da presente exposição é a ausência de «critérios
de equivalência» que permitam comparar o esquema de Picard de um esquema
projetivo com os de suas seções hiperplanas; os teoremas-chave para desenvolver
tais critérios encontram-se em [5], aos quais devem ser acrescentados os
teoremas de existência dos esquemas de Picard.

\fsection{1}{Propriedades topológicas dos pré-esquemas em grupos comutativos}

Sejam primeiro $k$ um corpo e $G$ um pré-esquema em grupos sobre $k$. Como o
elemento neutro $e$, por ser racional sobre $k$, é necessariamente fechado,
segue-se imediatamente que a diagonal de $G\times_kG$ é fechada e, portanto,
que $G$ é separado: todo pré-esquema em grupos sobre um corpo é separado.
Denotaremos por $G^\circ$ a componente conexa do elemento neutro $e$. Como $e$
é racional sobre $k$, $G^\circ$ é, de fato, geometricamente conexa, isto é, a
formação de $G^\circ$ é compatível com a extensão do corpo de base. Segue-se
também que $G^\circ$ é estável por multiplicação como conjunto e, quando $G$ é
localmente noetheriano, de modo que $G^\circ$ é aberto, pode-se considerar
$G^\circ$ como um subgrupo aberto de $G$. Doravante, suporemos que $G$ é
localmente de tipo finito sobre $k$; então $G^\circ$ é geometricamente
irredutível e de tipo finito sobre $k$. Com efeito, pode-se supor que $k$ é
algebricamente fechado e, além disso, que $G$ é reduzido (pois
$G_{\mathrm{red}}$ será então um subgrupo de $G$, tendo em vista que
$G_{\mathrm{red}}\times_kG_{\mathrm{red}}$ será ainda reduzido), logo liso
sobre $k$ em um aberto não vazio e, portanto, em toda parte, como se vê
transladando esse aberto. Mas então $G$ é localmente irredutível, de modo que
suas componentes irredutíveis coincidem com suas componentes conexas, e
$G^\circ$ é irredutível. Seja então $U$ uma vizinhança afim de $e$ em
$G^\circ$; usando a irredutibilidade de $G^\circ$, vê-se imediatamente que
$U\cdot U=G^\circ$, o que prova que $G^\circ$ é quase-compacto e, portanto, de
tipo finito sobre $k$.

Para simplificar, suponhamos que $G$ seja comutativo. Para todo inteiro $n>0$,
seja $G^{(n)}$ a imagem inversa de $G^\circ$ pelo homomorfismo $\varphi_n$ de
elevação à $n$-ésima potência em $G$; assim, $G^{(n)}$ é um subgrupo aberto de
$G$.
Ponhamos:
\[
G^\tau=\bigcup_nG^{(n)}.
\]
\src{223}
\[
G^\sigma=\bigcup_{(n,p)=1}G^{(n)},
\qquad
G^\rho=\bigcup_hG^{(p^h)},
\]
onde $p$ é o expoente característico do corpo $k$. Obtêm-se, portanto,
subgrupos abertos de $G$ que satisfazem
\[
G^\sigma\cap G^\rho=G^\circ,
\qquad
G^\sigma G^\rho=G^\tau.
\]

\result{OBSERVAÇÃO}{} Pode-se construir o esquema em grupos quociente
$G/G^\circ=N$ (cf. [3], IV) e definir então $G^\tau$, $G^\sigma$ e
$G^\rho$ como as imagens inversas em $G$ do subgrupo de torção de $N$
(respectivamente, de sua componente $p$-primária, respectivamente, do
complemento natural desta última, soma das componentes $q$-primárias para
$q$ primo $\neq p$).\footnote{O texto impresso troca as duas últimas
correspondências: atribui a $G^\sigma$ a componente $p$-primária e a
$G^\rho$ seu complemento primo com $p$. As definições precedentes de
$G^\sigma$ e $G^\rho$ impõem a ordem inversa; este ramo preserva a leitura
impressa.}
Observe-se, a esse respeito, que $N$ é um esquema em grupos discreto separável
sobre $k$ e, portanto, uma vez escolhida uma clausura algébrica $\overline k$
de $k$, que dá lugar a um grupo de Galois $\pi$, identifica-se com um grupo
discreto ordinário sobre o qual $\pi$ age por automorfismos. É dessa maneira
que se pode interpretar de modo evidente a construção do subgrupo de torção e
a decomposição deste em suas componentes $q$-primárias. Quando $G$ é o esquema
de Picard de um esquema $X$ próprio sobre $k$, $N$ pode ser chamado
\emph{esquema de Néron--Severi} (reduzido) de $X$ sobre $k$. Quando
$G^\circ_{\mathrm{red}}$ é um subesquema em grupos de $G^\circ$, o que
ocorre, em particular, sempre que $k$ é perfeito ou que $G^\circ$ é próprio
sobre $k$ (por exemplo, quando $X$ é geometricamente normal), convém introduzir
também o quociente $N'=G/G^\circ_{\mathrm{red}}$, que tende a comportar-se
melhor do que $N$ quanto à especialização, isto é, quando $X$ varia em uma
família de esquemas algébricos.

Sejam agora $S$ um pré-esquema localmente noetheriano e $G$ um pré-esquema em
grupos sobre $S$, localmente de tipo finito sobre $S$. Não supomos que $G$ seja
de tipo finito sobre $S$, nem separado sobre $S$. Poremos então
\[
G^\circ=\bigcup_{s\in S}(G_s)^\circ,
\]
e, quando $G$ for comutativo,
\[
G^\tau=\bigcup_{s\in S}(G_s)^\tau,
\qquad
G^\sigma=\bigcup_{s\in S}(G_s)^\sigma,
\qquad
G^\rho=\bigcup_{s\in S}(G_s)^\rho.
\]
\src{224}
Trata-se de subconjuntos de $G$, estáveis pela multiplicação em $G$, o que
evidentemente não implica que possam ser definidos mediante subpré-esquemas em
grupos de $G$. Em particular, parece não existir, em geral, um subpré-esquema
em grupos de $G$ cujo conjunto subjacente seja $G^\circ$. Naturalmente, se um
desses conjuntos for aberto, então, munido da estrutura induzida, será um
subpré-esquema aberto em grupos de $G$. Veremos que isso sempre ocorre com
$G^\tau$; assim, do ponto de vista dos funtores representáveis, e em particular
do ponto de vista das «especializações», a equivalência numérica comporta-se
de maneira mais satisfatória do que a equivalência algébrica. Eis as principais
propriedades gerais dos conjuntos que acabamos de definir:

\result{TEOREMA}{1.1} $G^\circ$, $G^\tau$, $G^\sigma$ e $G^\rho$ são
localmente construtíveis. Além disso:

\noindent (i) $G^\circ$ é quase-compacto sobre $S$. Quando os $G_s^\circ$ são
próprios e $G$ é separado sobre $S$, então $G^\circ$ é próprio sobre $S$ e,
portanto, fechado em $G$.

\noindent (ii) $G^\tau$ é aberto. Se $G^\circ$ é fechado, o mesmo vale para
$G^\tau$.

\noindent (iii) Se $G^\circ$ é fechado, o mesmo vale para $G^\sigma$, desde
que se esteja em característica igual, isto é, que todos os corpos residuais de
$S$ tenham a mesma característica. Se $G^\circ$ é fechado e
$G\longrightarrow S$ é universalmente aberto nos pontos de $G^\sigma$ (cf.
Corolário 1.5 abaixo), então $G^\sigma\longrightarrow S$ é universalmente
aberto.

\noindent (iv) Se $G^\circ$ é fechado, o mesmo vale para $G^\rho$. Suponhamos
que se esteja em característica igual e que, para todo inteiro $n>0$ tal que
$(n,p)=1$, o homomorfismo de elevação à $n$-ésima potência em $G$ seja aberto;
então $G^\rho$ é aberto.

Demos algumas indicações sobre a demonstração. O fato de $G^\circ$ ser
localmente construtível e quase-compacto sobre $S$ está contido no lema
seguinte:

\result{LEMA}{1.2} Suponhamos que $S\neq\varnothing$. Então existem um aberto
não vazio $U$ em $S$, um esquema em grupos $H$ de tipo finito sobre $U$ com
fibras conexas e, por fim, um homomorfismo de esquemas em grupos
$H\longrightarrow G|U$ que é uma imersão aberta cuja imagem é
$G^\circ|U$.

Para demonstrar o lema, pode-se supor que $S$ seja irredutível; seja $\eta$ seu
ponto genérico, e façamos a mudança de base
$S'=\operatorname{Spec}(\mathcal O_{S,\eta})\longrightarrow S$. Obtém-se um
esquema em grupos $G'$ sobre um anel artiniano local
$\mathcal O_{S,\eta}=A$, no qual há um subesquema aberto em grupos
$G'{}^\circ$ de tipo finito sobre $A$, como se disse acima. Este provém,
portanto, de um esquema em grupos $H$ de tipo finito sobre uma vizinhança
aberta $U$ de $\eta$, e a imersão canônica
$G'{}^\circ\longrightarrow G'$ provém de uma imersão
\src{225}
aberta $H\longrightarrow G|U$, que será um homomorfismo de esquemas em grupos
para $U$ suficientemente pequeno. Como as fibras de $H$ são conexas quando se
toma $U$ suficientemente pequeno, e têm todas a mesma dimensão, igual à das
fibras de $G$ para $U$ suficientemente pequeno, segue-se que, para todo
$s\in U$, a imagem de $H_s$ em $G_s$ é exatamente $G_s^\circ$ (após reduzir
novamente $U$), o que demonstra 1.2.

A segunda afirmação de (i) está contida no lema seguinte (que se aplica a uma
vizinhança aberta quase-compacta de $G^\circ$ em $G$):

\result{LEMA}{1.3} Sejam $X$ um pré-esquema separado de tipo finito sobre um
$S$ localmente noetheriano, $g$ uma seção de $X$ sobre $S$ e $X^\circ$ a união
das componentes conexas dos $g(s)$ nos $X_s$. Seja $s\in S$ tal que
$X_s^\circ$ seja próprio sobre $k(s)$. Então existe uma vizinhança aberta $U$
de $s$ tal que $X^\circ|U$ seja próprio sobre $U$ e, a fortiori, fechado em
$X|U$.

Por descida fielmente plana da base, reduzimo-nos ao caso em que $S$ é o
espectro de um anel local completo e $s$ é seu ponto fechado. Aplicando [2],
III, 5.5.1, vê-se que $X$ se decompõe na união disjunta de dois abertos $X'$ e
$X''$, o primeiro próprio sobre $S$ e tal que $X'_s=X_s^\circ$. Isso nos reduz
ao caso em que $X=X'$, isto é, em que $X$ é próprio sobre $S$. Nesse caso,
conclui-se por uma demonstração padrão que utiliza o critério valuativo para
que uma parte seja própria (omitido em [2], cap. II).

Demonstremos que $G^\tau$ é aberto ou, o que equivale ao mesmo, tendo em vista
que a formação de $G^\tau$ (assim como a de $G^\circ$, $G^\sigma$ e $G^\rho$)
comuta com a extensão da base: para toda seção $g$ de $G$ sobre $S$,
$g^{-1}(G^\tau)$ é aberto. Isso significa duas coisas:

\medskip
\noindent a. Seja $y\in S$ tal que $g(y)\in G_y^\tau$. Então, para todo
$y'\in\overline{\{y\}}$ suficientemente próximo de $y$, tem-se
$g(y')\in G_{y'}^\tau$.

\medskip
\noindent b. Seja $y'\in S$ tal que $g(y')\in G_{y'}^\tau$. Então, para toda
generalização $y$ de $y'$, tem-se $g(y)\in G_y^\tau$.

Para (a), observe-se que existe um inteiro $n>0$ tal que
$g^n(y)\in G_y^\circ$; como $G^\circ$ é construtível,
$(g^n)^{-1}(G^\circ)$ é construtível e, portanto, tem-se
$g^n(y')\in G_{y'}^\circ$ para $y'\in\overline{\{y\}}$ suficientemente
próximo de $y$. Para (b), observe-se que as potências
$g(y')^n=g^n(y')$ permanecem em um aberto quase-compacto de $G_{y'}$ (pois
estão contidas em um número finito de classes módulo $G_{y'}^\circ$); logo,
existe um aberto quase-compacto $U$ em $G$ que contém os $g^n(y')$ e, portanto,
também suas generalizações $g^n(y)$. Assim, as potências de $g(y)$ permanecem
em um aberto quase-compacto de $G_y$, o que implica facilmente que
$g(y)\in G_y^\tau$.
\src{226}
Suponhamos $G^\circ$ fechado; demonstremos que o mesmo vale para $G^\tau$.
Como já se sabe que $G^\tau$ é aberto e, portanto, localmente construtível,
resta demonstrar que é estável por especialização, o que decorre do fato de ser
uma união de fechados, a saber, as imagens inversas, pelos homomorfismos
$\varphi_n$ de elevação à $n$-ésima potência, do fechado $G^\circ$.

O mesmo argumento demonstrará que $G^\sigma$ e $G^\rho$ são fechados se
$G^\circ$ o for (sob a condição, no primeiro caso, de se estar em
característica igual), uma vez demonstrado que $G^\sigma$ e $G^\rho$ são
localmente construtíveis. Ora, para $x\in G^\tau$, seja $\nu(x)$ o menor
inteiro $n>0$ tal que o homomorfismo $\varphi_n$ de elevação à $n$-ésima
potência envie $x$ a $G^\circ$. Assim, $G^\sigma$ (respectivamente,
$G^\rho$) é formado pelos $x\in X$\footnote{O texto impresso traz aqui
«$X$», que não está definido nesta passagem; os quatro conjuntos considerados
são subconjuntos de $G$. Este ramo diplomático preserva a leitura impressa; o
ramo corrente usa $G$.} tais que $\nu(x)$ seja primo com $p$
(respectivamente, uma potência de $p$), e nossa afirmação de construtibilidade
decorre então da seguinte afirmação mais precisa:

\result{LEMA}{1.4} A função $\nu$ sobre $G^\tau$ é localmente construtível.

Isso significa, com efeito, que, para todo inteiro $n>0$, o conjunto dos
$x\in G^\tau$ tais que $\nu(x)=n$ é localmente construtível; ora, ele é a
diferença entre $\varphi_n^{-1}(G^\circ)$ e a união dos
$\varphi_d^{-1}(G^\circ)$, onde $d$ percorre os divisores próprios de $n$.
Como $G^\circ$ é localmente construtível, o mesmo vale para os
$\varphi_d^{-1}(G^\circ)$ e, portanto, também para a diferença anterior.

Suponhamos $G^\circ$ fechado e $G\longrightarrow S$ universalmente aberto nos
pontos de $G^\sigma$; demonstremos que $G^\sigma\longrightarrow S$ é aberto,
isto é, que transforma uma vizinhança em $G^\sigma$ de um $x\in G^\sigma$
em uma vizinhança de $y=f(x)$. Como $G^\sigma$ é localmente construtível,
reduzimo-nos a demonstrar que, para toda generalização $y'$ de $y$, existe uma
generalização $x'$ de $x$ em $G^\sigma$ sobre $y'$. Isso nos reduz, por
mudança de base, ao caso em que $S$ é o espectro de um anel de valorização
discreta e em que $y,y'$ são, respectivamente, o ponto fechado e o ponto
genérico. Usando o fato de que $G\longrightarrow S$ é aberto em $x$ (logo,
existe uma generalização $x_1$ de $x$ em $G$ sobre $y'$), pode-se supor ainda
que existe uma seção $g$ de $G$ sobre $S$ tal que $x=g(y)$ (após fazer outra
mudança de base). Se $k(y')$ tem característica 0, basta tomar qualquer
generalização $x'$ de $x$ em $G$ sobre $y'$; ela pertence a $G^\tau$ porque
$G^\tau$ é aberto e, portanto, a $G^\sigma$ porque
$G_{y'}^\sigma=G_{y'}^\tau$. Se a característica de $k(y')$ é $p>0$,
ponhamos
\[
\nu(g(y'))=p^hm\qquad\text{com }(m,p)=1 ;
\]
\src{227}
sejam $a$ e $b$ dois inteiros tais que $ap^h+bm=1$; ponhamos
\[
g_1=g^{ap^h},\qquad g_2=g^{bm},\qquad\text{logo}\qquad g=g_1g_2 ;
\]
portanto, por construção, tem-se $g_1(y')\in G^\sigma$ e
$g_2(y')\in G^\rho$. Como $G^\circ$ é fechado, segue-se que
$g_1(y)\in G^\sigma$ e $g_2(y)\in G^\rho$; daí, em virtude de
\[
g(y)=g_1(y)g_2(y)\in G^\sigma,
\]
também $g_2(y)\in G^\sigma$ e, portanto,
\[
g_2(y)\in G_y^\sigma\cap G_y^\rho=G_y^\circ.
\]
Ora, da hipótese e do fato de que $S$ é o espectro de um anel de valorização
discreta, resulta que $G-(G_y-G_y^\circ)$ é um aberto de $G$ sobre o qual
$G\longrightarrow S$ induz um morfismo aberto; logo, por todo ponto de
$G_y^\circ$ passa uma «quase-seção». Por conseguinte, após fazer ainda uma
extensão da base $S$, pode-se supor que existe uma seção $g'_2$ de $G^\circ$
sobre $S$ tal que $g'_2(y)=g_2(y)$. Ponhamos $g'=g_1g'_2$; então, por
construção, $g'(y)=g(y)=x$, e
\[
g'(y')=g_1(y')g'_2(y')\in G^\sigma,
\]
logo $g(y')=x'$\footnote{O texto impresso traz aqui «$g(y')=x'$»; a seção
construída na frase precedente e cuja fibra genérica está sendo calculada é
$g'$. Este ramo diplomático preserva a leitura impressa; o ramo corrente usa
$g'(y')=x'$.} é uma generalização de $x$ em $G^\sigma$ sobre $y'$, o que demonstra
(iii).

Demonstremos, por fim, a última afirmação (iv). Reduzimo-nos a provar que, se
$x\in G^\rho$, então toda generalização $x'$ de $x$ pertence a $G^\rho$.
Pode-se supor (após tomar as imagens por $\varphi_{n^h}$ para um $h$
conveniente)\footnote{O texto impresso traz $\varphi_{n^h}$; para
$x\in G^\rho$, a definição dada na p.~223 fornece uma potência do expoente
característico $p$ que envia $x$ a $G^\circ$. Este ramo diplomático preserva
a leitura impressa; o ramo corrente usa $\varphi_{p^h}$.} que se tem até mesmo
$x\in G^\circ$. Então, para todo inteiro $n$ primo com a característica, $x$
pertence à imagem de $\varphi_n$ (pois a elevação à $n$-ésima potência em um
grupo conexo de tipo finito sobre um corpo de característica prima com $n$ é
sobrejetiva). Como $\varphi_n$ é aberta, segue-se que $x'$ também pertence à
imagem de $\varphi_n$. Mais precisamente, seja $U$ um aberto quase-compacto de
$G^\tau$ que contém $G_y^\circ$; então se tem $x'\in\varphi_n(U)$ para todo
$n$ primo com $p$. Tomando para $n$ um múltiplo comum dos fatores primos com
$p$ dos $\nu(z)$, para $z\in U$, obtém-se $x'\in G^\rho$.

A afirmação (1.1), (iii), completa-se assim:

\result{COROLÁRIO}{1.5} Seja $n>1$ um inteiro tal que o homomorfismo de
elevação à $n$-ésima potência $\varphi_n:G\longrightarrow G$ seja
universalmente aberto (por exemplo, étale); ponhamos
\[
G^{(n)}=\varphi_n^{-1}(G^\circ),
\]
e suponhamos que as fibras conexas $G_s^\circ$ não «contenham
\src{228}
componente aditiva» (isto é, o grupo obtido pela extensão do corpo $k(s)$ a
uma clausura algébrica não contém nenhum subgrupo isomorfo a
$\mathbf G_a$). Então $G\longrightarrow S$ é universalmente aberto nos
pontos dos $G^{(n^h)}$. Em particular, se os $G_s^\circ$ não contêm componente
aditiva e se, para todo inteiro $n>1$, o homomorfismo $\varphi_n$ é
universalmente aberto nos pontos de característica residual prima com
$p$\footnote{O texto impresso traz «prima com $p$». Nesta proposição
quantificada por $n$, o critério de etalidade e o argumento pelos núcleos
${}_nG$ exigem característica residual prima com $n$. Este ramo diplomático
preserva $p$; o ramo corrente usa $n$.}, então
$G\longrightarrow S$ é universalmente aberto nos pontos de $G^\sigma$ e,
portanto, por (1.1), (iii), se $G^\circ$ é fechado,
$G^\sigma\longrightarrow S$ é universalmente aberto. Nesses casos,
$s\longmapsto\dim G_s$ é uma função localmente constante sobre $S$.

Com efeito, da hipótese resulta que o núcleo ${}_nG$ de $\varphi_n$ é
universalmente aberto sobre $S$; logo, $G\longrightarrow S$ é universalmente
aberto nos pontos de ${}_nG$ e, portanto (substituindo $n$ por $n^h$), nos
pontos dos ${}_{n^h}G$. Ora, a hipótese sobre as fibras $G_s^\circ$ serve
exatamente para assegurar que os pontos de ordem uma potência de $n$ em $G_s$
são densos na união dos $G^{(n^h)}$; daí resulta facilmente que
$G\longrightarrow S$ é universalmente aberto em todos os pontos dessa união,
em particular ao longo da seção neutra, de onde decorre facilmente que
$s\longmapsto\dim G_s$ é uma função localmente constante.

\result{OBSERVAÇÃO}{1.6} Recordemos que um morfismo $X\longrightarrow S$ é
dito universalmente aberto se transforma todo aberto em um aberto e conserva
essa louvável propriedade após toda mudança de base $S'\longrightarrow S$.
Isso significa também (quando $S$ é localmente noetheriano e
$X\longrightarrow S$ é localmente de tipo finito) que toda componente
irredutível de $X$ domina $S$ e que essa propriedade é preservada após toda
mudança de base $S'\longrightarrow S$. Em ambas as afirmações, basta, aliás
(sob as hipóteses de finitude acima), testar mudanças de base
$S'\longrightarrow S$, em que $S'$ é o espectro de um anel de valorização
discreta (completo, com corpo residual algebricamente fechado, se assim se
quiser~...). A definição estende-se de maneira evidente ao caso de um
subconjunto $Z$ de $X$ (como o subconjunto $G^\sigma$ de $G$). É um fenômeno
perfeitamente normal, mesmo partindo de um morfismo projetivo liso
$X\longrightarrow S$ com fibras geométricas conexas (por exemplo, o produto
fibrado da família modular de curvas elípticas consigo mesma sobre
$S=\operatorname{Spec}\mathbf C[j]$), que $\operatorname{Pic}_{X/S}$ não
seja universalmente aberto sobre $S$, isto é, que possa haver componentes
irredutíveis de $\operatorname{Pic}_{X/S}$ situadas inteiramente sobre um
único ponto de $S$: isso está ligado ao fato de que o posto do grupo de
Néron--Severi das fibras de $X/S$ pode apresentar saltos para cima (fenômenos
de «multiplicação complexa»). Em contrapartida, (1.5) assegura-nos que, nos
casos favoráveis, $\operatorname{Pic}_{X/S}^\sigma$ (e, ao que parece, na
maioria das vezes até mesmo $\operatorname{Pic}_{X/S}^\tau$) é
universalmente aberto sobre $S$.
\src{229}
Por fim, eis um caso útil em que, excepcionalmente, $G^\circ$ consente em ser
aberto:

\result{COROLÁRIO}{1.7} Suponhamos que $G\longrightarrow S$ seja
universalmente aberto nos pontos de $G^\circ$ (cf. (1.5)) e que as fibras
$G_s$ sejam separáveis e, portanto, lisas sobre $k(s)$ (esta última condição é automaticamente satisfeita em
característica residual zero, em virtude de um resultado de Cartier). Então
$G^\circ$ é aberto em $G$. Se, além disso, $S$ é reduzido, $G^\circ$ é liso
(em particular, plano) sobre $S$.

A primeira afirmação pode, aliás, ser precisada observando-se que, se, para um
$s\in S$ dado, $G\longrightarrow S$ é universalmente aberto nos pontos de
$G_s^\circ$ e $G_s^\circ$ é separável sobre $k(s)$, então $G^\circ$ é uma
vizinhança de $G_s^\circ$ (e, aliás, a hipótese feita em $s$ continuará
verificada nos pontos vizinhos). Esse enunciado, independente de toda estrutura
de grupos, será demonstrado em [2], IV, §~7. A última afirmação de (1.7), também
independente de toda estrutura de grupos e de toda questão de componente
conexa, é um caso particular de um critério de planicidade dado em [IV], §~5,
que implica, mais geralmente:

\result{COROLÁRIO}{1.8} Seja $U$ um subconjunto aberto de uma fibra $G_s$ tal
que $G\longrightarrow S$ seja universalmente aberto nos pontos de $U$ (cf.
(1.5)). Se $G_s$ é separável sobre $k(s)$ e $S$ é reduzido em $s$, então $G$ é
plano sobre $S$ nos pontos de $U$ (e, portanto, neste caso, liso sobre $S$ nos
pontos de $U$).

\result{OBSERVAÇÕES}{1.9} Ignoro se, quando $G$ é separado sobre $S$,
$G^\circ$ ou $G^\tau$ é sempre fechado em $G$; isso parece pouco provável.
Encontram-se, em todo caso, contraexemplos evidentes se se abandona a hipótese
de separação, por exemplo tomando a reta afim cuja origem é duplicada uma
infinidade enumerável de vezes, obtendo-se assim um esquema em grupos $G$ sobre
a reta afim cujas fibras se reduzem todas ao grupo unidade, salvo uma cuja
fibra é $\mathbf Z$. Esse esquema em grupos é um subesquema aberto em grupos,
a saber, o fecho da seção unidade, no pré-esquema de Picard do $S$-esquema $X$
correspondente a uma família de cônicas que degenera em duas retas concorrentes.

Além disso, mesmo partindo de um esquema em grupos finito e plano $G$ sobre o
espectro $S$ de um anel de valorização discreta $V$ e, por conseguinte, se
$G^\circ$ se reduz à seção unidade, sendo, portanto, fechado, várias
conclusões tornam-se falsas quando se abandonam certas hipóteses. Suponhamos
que $V$ seja de característica igual $p>0$ e seja $G$ o núcleo do homomorfismo
$\mathbf G_a\longrightarrow\mathbf G_a$ definido pelo homomorfismo de
funtores
\[
f\longmapsto f^p-tf,
\]
onde $t$ é um uniformizante de $V$.
\src{230}
(Lembremos que, por definição, o «grupo aditivo» $\mathbf G_a$ sobre $S$
representa o funtor $S'\longmapsto\Gamma(S',\mathcal O_{S'})$.) Então, como
$S$-esquema,
\[
G=\operatorname{Spec}V[x]/(x^p-tx)
\]
é a união de $p$ «retas» concorrentes, representando um grupo
$\mathbf Z/p\mathbf Z$ que degenera em um grupo infinitesimal de tipo aditivo.
Neste exemplo, $G^\circ=G^\sigma$ (e é uma das $p$ retas), portanto não é
aberto em $G$. No caso de característica mista, de característica residual
$p>0$, partamos do esquema em grupos $H=\boldsymbol\mu_p$, núcleo da aplicação
de elevação à $p$-ésima potência em $\mathbf G_m$, portanto
\[
H=\operatorname{Spec}V[x]/(x^p-1) ;
\]
é ainda a união de $p$ retas concorrentes, representando um grupo
$\mathbf Z/p\mathbf Z$ (em característica zero) que degenera em um grupo
infinitesimal de tipo multiplicativo em característica $p$. Seja $H'$ o
«esquema em grupos constante» definido pelo grupo finito ordinário
$\mathbf Z/p\mathbf Z$, portanto a soma disjunta, indexada por
$\mathbf Z/p\mathbf Z$, de cópias de $S$, ou ainda
$\operatorname{Spec}V^{\mathbf Z/p\mathbf Z}$. Então $G=H\times_S H'$ descreve
um grupo $(\mathbf Z/p\mathbf Z)^2$ em característica zero, que degenera em um
grupo infinitesimal vezes $\mathbf Z/p\mathbf Z$ em característica $p$. Aqui
$G^\rho$ é a união da seção unidade e da fibra especial, portanto $G^\rho$ não
é aberto, ao passo que $G^\sigma$ é a união da seção unidade e da fibra
genérica, portanto não é fechado, contrariamente ao que se afirma nos casos de
característica igual em (1.1), (iii) e (iv). Naturalmente, esses fenômenos estão
ligados à característica $p>0$. Os resultados precedentes dão, com efeito:

\result{COROLÁRIO}{1.10} Suponhamos que as características residuais de $S$
sejam todas nulas, portanto $G^\tau=G^\sigma$ e $G^\circ=G^\rho$. Então
$G^\tau=G^\sigma$ é aberto, e até aberto e fechado se $G$ é separado sobre $S$
e se os $G_s^\circ$ são próprios; sob essa mesma hipótese, $G^\circ=G^\rho$ é
próprio sobre $S$, portanto fechado em $G$. Suponhamos, por fim, que, para todo
inteiro $n>1$, o homomorfismo de elevação à $n$-ésima potência em $G$ seja
universalmente aberto; então $G^\circ$ é aberto; e, se além disso os
$G_s^\circ$ não têm componente aditiva (por exemplo, se os $G_s^\circ$ são
próprios, como acima), então $G^\tau\longrightarrow S$ é universalmente aberto,
e até liso se $S$ é reduzido.

Assinalemos, por fim, o seguinte resultado fácil:

\result{PROPOSIÇÃO}{1.11} Existe uma parte aberta $U$ de $S$ tal que o conjunto
$G'$ dos pontos de $G$ nos quais $G$ é liso (resp. plano) sobre $S$ seja o
conjunto subjacente a um subesquema em grupos aberto induzido de $G|U$. Além
disso, toda seção de $G$ sobre $U$ é uma seção de $G'$ sobre $U$.
\src{231}
\result{COROLÁRIO}{1.12} Se $G$ é liso (resp. plano) sobre $S$ nos pontos da
seção unidade, também o é nos pontos de toda seção de $G$
sobre $S$ e em todos os pontos de $G^\circ$. Se, além disso, para todo inteiro
$n>0$, o homomorfismo de elevação à $n$-ésima potência
$\varphi_n:G\longrightarrow G$ é étale nos pontos de característica prima com
$n$, então $G$ é liso (resp. plano) sobre $S$ em todos os pontos de $G^\sigma$.

\fsection{2}{Aplicação às propriedades locais dos esquemas de Picard}

\result{TEOREMA}{2.1}

\noindent(i) Seja $f:X\longrightarrow S$ um morfismo próprio e liso, e
suponhamos que $\operatorname{Pic}_{X/S}$ exista (por exemplo, que $f$ seja um
morfismo projetivo). Então $\operatorname{Pic}_{X/S}$ é separado sobre $S$ e,
para toda parte fechada $Z$ de $\operatorname{Pic}_{X/S}$ que seja de tipo
finito sobre $S$, $Z$ é próprio sobre $S$.

\noindent(ii) Seja $X$ um pré-esquema sobre um corpo $k$, próprio e
geometricamente normal. Então $\operatorname{Pic}_{X/S}^\circ$\footnote{O ramo
corrente traz $\operatorname{Pic}_{X/k}^\circ$, pois $S$ não intervém neste
enunciado sobre o corpo $k$.} é próprio sobre $k$.

\begin{proof}
(i) Os critérios valuativos ([2], II, parágrafo 7) reduzem-nos a provar o
seguinte: se $S$ é o espectro de um anel de valorização discreta completo e $U$
o aberto reduzido ao ponto genérico de $S$, então toda seção racional de
$\operatorname{Pic}_{X/S}$ sobre $S$, isto é, toda seção sobre $U$, prolonga-se
de maneira única a uma seção sobre $S$. Tendo em conta a definição de
$\operatorname{Pic}_{X/S}$, isso equivale ao seguinte enunciado: para todo
módulo invertível $\mathcal L$ sobre $V=f^{-1}(U)$, existe um módulo invertível
sobre $X$ que prolonga $\mathcal L$, único a menos de isomorfismo. Ora, isso
decorre facilmente da descrição dos módulos invertíveis sobre $V$, resp. $X$,
em termos de classes de divisores de Cartier, tendo em conta que os anéis
locais de $X$ são regulares ($X$ sendo liso sobre $S$ regular), portanto
fatoriais pelo teorema de Auslander, o que implica que todo divisor sobre
$S$\footnote{O ramo corrente traz aqui «sobre $X$»; a fatorialidade invocada é
a dos anéis locais de $X$.} é um divisor de Cartier. Com efeito, todo divisor
sobre $V$ pode ser prolongado a um divisor sobre $X$ tomando-se o seu fecho.
\end{proof}

\result{OBSERVAÇÃO}{2.2} A demonstração permanece válida supondo-se apenas que
$f$ é plano e que suas fibras $X_s$ são localmente interseções completas e
lisas sobre $k(s)$ em codimensão $\leq2$, tendo em conta o seguinte resultado
demonstrado em [5]: um anel local noetheriano de interseção completa que é
regular em codimensão $\leq3$ é fatorial («conjectura de Samuel»). Note-se que o
\src{232}
resultado torna-se falso se se substitui «codimensão $\leq2$» por «codimensão
$\leq1$», isto é, pela hipótese «normal», como se pode constatar no exemplo de
uma família de quádricas não singulares que degenera em um cone quadrático.

(ii) Utilizando o lema de Chow, reduzimo-nos ao caso em que $X$ é projetivo,
portanto imerso em $\mathbf P_k^n$; podemos supor $X$ conexo. Se $\dim X=1$,
então $X$ é liso sobre $k$, e aplica-se (i). Se
$\dim X\geq2$, pode-se utilizar uma variante dos «critérios de equivalência»
conhecidos, que implica existir um número finito de curvas $Y_i$ lisas em $X$
(obtidas como interseções de $X$ com subespaços lineares convenientes de
$\mathbf P_k^n$), tais que
\[
\operatorname{Pic}_{X/k}^\tau\longrightarrow
\prod_i\operatorname{Pic}_{Y_i/k}^\tau
\]
seja um monomorfismo e induza, a fortiori, um monomorfismo nas componentes
conexas. Como a componente conexa do segundo membro é própria sobre $k$ pelo
que precede, e como se trata de um homomorfismo de esquemas em grupos, que será
necessariamente uma imersão fechada, segue-se que
$\operatorname{Pic}_{X/k}^\circ$ também é próprio sobre $k$. Pode-se evitar o
recurso aos delicados critérios de equivalência utilizando-se a estrutura dos
grupos algébricos comutativos sobre um corpo algebricamente fechado (devida a
Chevalley--Borel); reduzimo-nos a provar que todo morfismo da reta afim
\footnote{A leitura impressa «da reta afim» é preservada neste ramo; a errata
vinculada, p.~303, e o ramo corrente prescrevem «reta afim privada da origem».}
em $\operatorname{Pic}_{X/k}^\tau$ é constante, o que equivale a dizer que todo
módulo invertível sobre $X[t]$\footnote{A leitura impressa é preservada neste
ramo; a errata vinculada, p.~303, e o ramo corrente usam $X[t,t^{-1}]$ em lugar
de $X[t]$.} provém de um módulo invertível sobre $X$, resultado bem conhecido e
muito elementar, que nem sequer utiliza o fato de $X$ ser próprio sobre $k$ (a
hipótese de $X$ ser normal permitindo reduzir-nos imediatamente ao caso em que
$X$ é regular).

\result{COROLÁRIO}{2.3} Seja $f:X\longrightarrow S$ um morfismo próprio e
normal (isto é, plano e com fibras geométricas normais), e suponhamos que
$\operatorname{Pic}_{X/S}$ exista; então $\operatorname{Pic}_{X/S}^\circ$ é
próprio sobre $S$, portanto fechado em $\operatorname{Pic}_{X/S}$; além disso,
$\operatorname{Pic}_{X/S}^\tau$ e $\operatorname{Pic}_{X/S}^\rho$ também são
fechados, assim como $\operatorname{Pic}_{X/S}^\sigma$ em «característica
igual».

Aplicam-se (1.1) e (2.1), (ii).

\result{COROLÁRIO}{2.4} Seja $f:X\longrightarrow S$ um morfismo próprio e liso,
tal que $\operatorname{Pic}_{X/S}$ exista e seja uma soma disjunta de esquemas
$P^{(i)}$ de tipo finito sobre $S$ (cf. [3], V, 4.1). Então cada $P^{(i)}$ é
próprio sobre $S$.

Isso decorre de (2.1), (i). Como se assinalou em (2.2), o resultado pode ser
generalizado fazendo-se hipóteses menos restritivas sobre as fibras de $f$, mas
torna-se falso se nos limitamos a supor $f$ normal. Nesse caso, ignoro, aliás,
se $\operatorname{Pic}_{X/S}^\tau$ é, não obstante, próprio sobre $S$, mesmo
admitindo que seja de tipo finito sobre $S$.
\src{233}
\result{TEOREMA}{2.5} Seja $f:X\longrightarrow S$ um morfismo próprio e plano,
tal que $\operatorname{Pic}_{X/S}$ exista, e, para todo inteiro $n$, seja
$\varphi_n$ o homomorfismo de elevação à $n$-ésima potência nesse pré-esquema
em grupos. Então $\varphi_n$ é étale em todo ponto
$x\in X$\footnote{O ramo corrente traz $x\in\operatorname{Pic}_{X/S}$, pois
$\varphi_n$ é o endomorfismo de potência do pré-esquema de Picard.} de
característica residual prima com $n$.

Pela caracterização infinitesimal dos morfismos étales, essa afirmação equivale
à seguinte:

\result{LEMA}{2.6} Suponhamos que $S$ seja o espectro de um anel artiniano local
$A$ cujo ideal maximal $\mathfrak m$ tenha potência $(\nu+1)$-ésima nula, e
sejam
\[
A_{\nu-1}=A/\mathfrak m^\nu,
\qquad X_{\nu-1}=X\otimes_A A_{\nu-1},
\]
$\mathcal L$ um módulo invertível sobre $X$ e $\mathcal L'_{\nu-1}$ um módulo
invertível sobre $X_{\nu-1}$ cuja potência tensorial $n$-ésima seja isomorfa a
\[
\mathcal L_{\nu-1}=\mathcal L\otimes_A A_{\nu-1}.
\]
Então existe um módulo invertível $\mathcal L'$ sobre $X$ cuja potência
tensorial $n$-ésima é isomorfa a $\mathcal L$ (se $n$ é primo com a
característica residual de $k=k(A)$).

Demonstremos o lema. Ponhamos
\[
V=\mathfrak m^\nu=\mathfrak m^\nu/\mathfrak m^{\nu+1} ;
\]
este é um espaço vetorial sobre $k=k(A)$. Começamos por prolongar
$\mathcal L'_{\nu-1}$ a um módulo invertível qualquer $\mathcal L'$ sobre $X$.
A obstrução a fazê-lo encontra-se em
\[
H^2(X_0,\mathcal O_{X_0})\otimes_k V,
\]
mas, em virtude da hipótese feita sobre $\mathcal L'_{\nu-1}$ e do fato de
$\mathcal L_{\nu-1}$ poder ser prolongado, vê-se que o produto dessa obstrução
por $n$ é nulo, portanto ela própria é nula, pois $n$ é primo com a
característica. Além disso, a arbitrariedade na extensão efetuada encontra-se
em
\[
H^1(X_0,\mathcal O_{X_0})\otimes_k V,
\]
e o desvio $\xi$ de $\mathcal L'^{\otimes n}$ em relação a $\mathcal L$
encontra-se no mesmo módulo; se se procura corrigir $\mathcal L'$ de modo a
tornar nulo esse desvio, reduzimo-nos a encontrar um $\eta$ no referido módulo
tal que $n\eta=\xi$. Isso é novamente possível porque $n$ é primo com a
característica.

\result{COROLÁRIO}{} Sob as condições de (2.5), suponhamos, além disso, que os
esquemas de Picard das fibras $X_s$ não contenham componente aditiva (por
exemplo, que os $X_s$ sejam geometricamente normais, cf. (2.1), (ii)). Então
$\operatorname{Pic}_{X/S}\longrightarrow S$ é universalmente aberto nos pontos
de $\operatorname{Pic}_{X/S}^\circ$. Se
$\operatorname{Pic}_{X/S}^\circ$ é fechado (por exemplo, se os $X_s$ são
geometricamente normais, cf. (2.3)), então
$\operatorname{Pic}_{X/S}^\sigma$ é ele próprio universalmente aberto sobre
$S$. Por fim, no caso de característica igual,
$\operatorname{Pic}_{X/S}^\rho\longrightarrow S$ é universalmente aberto.

Aplicam-se (1.5) e (1.1).
\src{234}
\result{COROLÁRIO}{2.7} Seja $f:X\longrightarrow Y$\footnote{O ramo corrente
traz $f:X\longrightarrow S$, pois tanto o pré-esquema de Picard quanto a função
seguinte têm base $S$.} um morfismo próprio e plano, tal que
$\operatorname{Pic}_{X/S}$ exista. Então a função
\[
s\longmapsto\dim\operatorname{Pic}_{X_s/k(s)}
\]
sobre $S$ é semicontínua superiormente (isto é, pode dar saltos para cima, mas
não para baixo), e é até contínua (isto é, localmente constante) se os
$\operatorname{Pic}_{X_s/k(s)}$ não contêm componente aditiva.

A primeira afirmação é trivial, ou quase, para todo pré-esquema em grupos
localmente de tipo finito sobre uma base localmente noetheriana, pois basta
examinar ao longo da seção unidade. A segunda afirmação decorre de (2.5).

\result{OBSERVAÇÃO}{2.8} Sejam $s,s'\in S$ tais que $s$ seja uma especialização
de $s'$. Então (2.7) equivale a uma desigualdade (resp. igualdade) entre as
dimensões dos esquemas de Picard de $X_{s'}$ e de sua «especialização» $X_s$.
Aliás, Serre havia observado, ainda antes da construção dos esquemas de Picard,
que a invariância da dimensão dos Picard dos $X_s$, no caso de um morfismo liso
$f:X\longrightarrow S$, era uma consequência formal da teoria da
especialização do grupo fundamental ([4], X) e das relações clássicas de Kummer
entre os pontos de ordem finita na variedade de Picard e o grupo fundamental
abelianizado ([4], XI). Se designamos, de modo geral, por $\alpha,\mu,\lambda$
as dimensões das partes abeliana, multiplicativa e aditiva de
$\operatorname{Pic}_{X_s/k(s)}^\circ$, e definimos do mesmo modo
$\alpha',\mu',\lambda'$, as relações conhecidas exprimem-se pelas seguintes
desigualdades:
\[
\tag{*}\alpha+\mu+\lambda\geq\alpha'+\mu'+\lambda'
\]
(válida contanto que $\operatorname{Pic}_{X/S}$ exista, portanto provavelmente
em todos os casos), desigualdade que, para $\lambda=\lambda'=0$, se reduz a uma
igualdade, válida sob as mesmas hipóteses de existência:
\[
\alpha+\mu=\alpha'+\mu' ;
\]
por outro lado, tem-se
\[
\tag{**}2\alpha+\mu\leq2\alpha'+\mu'
\]
\src{235}
pelo argumento de Serre, se os $X_s$ são separáveis (sem sequer supor a
existência de $\operatorname{Pic}_{X/S}^\tau$), ou se os
${}_n\operatorname{Pic}_{X/S}$ (núcleos de $\varphi_n$ no funtor de Picard) são
separados sobre $S$ (tendo em conta que são étales sobre $S$ graças a (2.5)).
Tender-se-ia a conjecturar que (*) é uma igualdade em todos os casos, ou pelo
menos se os $X_s$ são separáveis, e que se têm as desigualdades
\[
\tag{***}\alpha\leq\alpha',\qquad\lambda\geq\lambda',
\]
que deveriam ser válidas sempre que se tenha um pré-esquema em grupos,
localmente de tipo finito sobre $S$ localmente noetheriano, no qual a dimensão
das fibras seja constante (ver (1.3) para um resultado positivo nessa direção).

\result{OBSERVAÇÃO}{2.9} Em todos os casos conhecidos,
$\operatorname{Pic}_{X/S}^\tau$ é universalmente aberto sobre $S$, mas sem
dúvida não se deve tirar daí ilusões excessivas, mesmo que
$f:X\longrightarrow S$ seja liso; em todo caso, Mumford
construiu um exemplo (é verdade que com $S$ não reduzido, sendo de fato $S$ o
espectro de um anel artiniano) no qual $\operatorname{Pic}_{X/S}^\tau$ não é
plano sobre $S$, fazendo variar infinitesimalmente a superfície de Igusa. O
ponto considerado por Mumford encontra-se, aliás, em
$\operatorname{Pic}_{X/S}^\rho$, e continua possível (para
$f:X\longrightarrow S$ liso) que $\operatorname{Pic}_{X/S}^\tau$ seja plano
sobre $S$ nos pontos de $\operatorname{Pic}_{X/S}^\sigma$; o conferencista
duvida, contudo, que isso seja sempre assim, mesmo restringindo-se aos pontos de
$\operatorname{Pic}_{X/S}^\circ$ e supondo que $S$ seja o espectro de um anel
de valorização discreta. A questão está, aliás, ligada ao estudo dos pontos
fixos de um esquema abeliano sob um grupo finito de automorfismos, situação
para a qual parecem faltar exemplos.

Parece que, mesmo limitando-nos a $f:X\longrightarrow S$ liso e projetivo, os
resultados de regularidade local sobre $\operatorname{Pic}_{X/S}$ enunciados no
presente número e as conjecturas assinaladas em (2.8) esgotam aproximadamente o
que se pode dizer a esse respeito sem hipóteses mais particulares sobre a
natureza das fibras de $f$. Lembremos, contudo, que, se as fibras geométricas
de $\operatorname{Pic}_{X/S}$ são reduzidas e sem componente aditiva, então
decorre de (1.8) e (2.5) que $\operatorname{Pic}_{X/S}$ é liso sobre $S$ nos
pontos de $\operatorname{Pic}_{X/S}^\sigma$ quando $S$ é reduzido; esse
resultado vale até mesmo, se $f:X\longrightarrow S$ é normal, sem hipótese
sobre $S$, como veremos em (3.5). Assinalemos, a esse respeito:

\result{PROPOSIÇÃO}{2.10}

\noindent(i) Se $\operatorname{Pic}_{X/S}$ é liso (resp. plano) sobre $S$ nos
pontos da seção unidade, também o é em todos os pontos de
$\operatorname{Pic}_{X/S}^\sigma$ e nos pontos de toda seção de
$\operatorname{Pic}_{X/S}$ sobre $S$.
\src{236}
\noindent(ii) Seja $s\in S$ tal que
\[
H^2(X_s,\mathcal O_{X_s})=0,
\]
então existe uma vizinhança aberta $U$ de $s$ tal que
$\operatorname{Pic}_{X/S}|U$ seja liso sobre $U$.

\noindent(iii) Seja $X$ um esquema próprio sobre um corpo; então tem-se
\[
\dim\operatorname{Pic}_{X/k}\leq\dim H^1(X,\mathcal O_X),
\]
com igualdade se e somente se $\operatorname{Pic}_{X/k}$ é liso sobre $k$;
isso ocorre sempre se $k$ tem característica zero.

(i) decorre de (2.5) e (1.12), (ii) do critério infinitesimal de lisura e de um
cálculo de obstruções bem conhecido, tendo em conta que (pelo «teorema de
semicontinuidade») a hipótese feita em $s$ continuará verificada nos pontos
vizinhos. Por fim, (iii) decorre do fato de que $H^1(X,\mathcal O_X)$ é
isomorfo ao espaço tangente de Zariski no elemento neutro de
$\operatorname{Pic}_{X/k}$; a última afirmação é um caso particular de um
teorema de Cartier segundo o qual um «grupo formal» em característica zero é
formalmente liso sobre $k$.

\fsection{3}{O subesquema abeliano canônico de
$\operatorname{Pic}_{X/S}$ e o esquema de Albanese}

\result{PROPOSIÇÃO}{3.1} Sejam $k$ um corpo e $G$ um esquema em grupos de tipo
finito sobre $k$, comutativo e «sem componente aditiva». Então
$G_{\mathrm{red}}^\circ$ é separável sobre $k$ e, por conseguinte, é um
subesquema em grupos liso sobre $k$.

Sendo isso trivial se $k$ é perfeito, em particular para $G_{\bar k}$, onde
$\bar k$ é o fecho algébrico de $k$, basta ver que
$(G_{\bar k}^\circ)_{\mathrm{red}}$ provém de um subesquema de $G$. Ora, a
hipótese sobre $G_{\bar k}$ (ausência de componente aditiva)\footnote{A leitura
impressa equivalente a «Ora, a hipótese ... resulta» é preservada neste ramo;
a errata vinculada, p.~303, e o ramo corrente inserem a preposição necessária
para «Ora, da hipótese ... resulta».} resulta facilmente que existe um inteiro
$m$ tal que $(G_{\bar k}^\circ)_{\mathrm{red}}$ seja a imagem esquemática do
homomorfismo de elevação à $m$-ésima potência em $G_{\bar k}$. Como esse último
homomorfismo provém do homomorfismo análogo para $G$, a imagem esquemática deste
último responde à questão.

\result{COROLÁRIO}{3.2} Seja $X$ um esquema normal e próprio sobre $k$ tal que
$\operatorname{Pic}_{X/k}$ exista; então existe um subesquema abeliano $A$ de
$\operatorname{Pic}_{X/k}$ cujo conjunto subjacente é
$\operatorname{Pic}_{X/k}^\circ$.

Com efeito, em virtude de (2.1), (ii), sendo
$\operatorname{Pic}_{X/k}^\circ$ próprio sobre $k$, ele satisfaz a condição de
(3.1).
\src{237}
O resultado precedente mostra, portanto, que, em certos casos, a clássica
«variedade de Picard» (que, na teoria atual, seria denotada por
$(\operatorname{Pic}_{X/k})_{\mathrm{red}}^\circ$) «é definida sobre $k$»,
sem supor que o corpo $k$ seja perfeito.

Seja agora $f:X\longrightarrow S$ um esquema relativo próprio e plano, com
\[
\mathcal O_S\xrightarrow{\sim}f_*(\mathcal O_X)
\]
para simplificar, tal que $\operatorname{Pic}_{X/S}$ exista e que
$\operatorname{Pic}_{X/S}^\circ$ seja próprio sobre $S$. Suporemos, além
disso, para (3.3), (ii), que existe um aberto de $\operatorname{Pic}_{X/S}$
contendo $\operatorname{Pic}_{X/S}^\circ$ que seja quase-projetivo sobre $S$;
essa condição é satisfeita, como vimos, se $f$ é projetivo e tem fibras
geométricas separáveis e irredutíveis. Recordemos que se chama esquema
abeliano sobre $S$ um esquema em grupos sobre $S$, próprio e liso sobre
$S$, com fibras geométricas conexas. Propomo-nos
examinar se existe um subesquema em grupos $A$ de
$\operatorname{Pic}_{X/S}$ que seja um esquema abeliano e cujo conjunto
subjacente seja $\operatorname{Pic}_{X/S}^\circ$. Acabamos de ver que sempre
existe um se $S$ é o espectro de um corpo. Eis o que se sabe dizer no caso
geral considerado aqui:

\result{TEOREMA}{3.3} Sob as condições precedentes:

\noindent(i) Se existe um subesquema abeliano de
$\operatorname{Pic}_{X/S}$ cujo conjunto subjacente é
$\operatorname{Pic}_{X/S}^\circ$, ele é único. Sua formação é, por
conseguinte, compatível com as mudanças de base.

\noindent(ii) Para que exista tal subesquema abeliano, basta que exista um
após toda mudança de base $S'\longrightarrow S$, em que $S'$ é artiniano
local; se $S$ é o espectro de um anel local, basta até mesmo testar com os
$S'=\operatorname{Spec}(A_n)$, em que $A_n=A/\mathfrak m^{n+1}$. Se $S$ é
reduzido, basta igualmente testar com as mudanças de base
$S'\longrightarrow S$, em que $S'$ é o espectro de um anel de valorização
discreta (completo com corpo residual algebricamente fechado, se assim se
desejar).

\noindent(iii) Suponhamos que $A$ exista e seja
$B=\operatorname{Alb}^0(X/S)$ o esquema abeliano dual (isto é,
$B=\operatorname{Pic}_{A/S}^\circ$ [7]). Pode-se então construir
canonicamente um espaço homogêneo principal
$P=\operatorname{Alb}^1(X/S)$ sob $B$ e um $S$-morfismo
$X\longrightarrow P$ que seja universal para os $S$-morfismos de $X$ em
esquemas para-abelianos (isto é, espaços homogêneos principais sob esquemas
abelianos). A formação de $\operatorname{Alb}^0(X/S)$,
$\operatorname{Alb}^1(X/S)$ e
$X\longrightarrow\operatorname{Alb}^1(X/S)$ comuta com a mudança de base.

Esbocemos a demonstração:

(i) é uma propriedade geral de rigidez dos subesquemas abelianos dos esquemas
em grupos comutativos: se dois desses subesquemas coincidem como conjuntos em
um ponto $s\in S$, eles coincidem sobre toda a componente
\src{238}
conexa de $s$ [7]. (Este resultado generaliza um teorema clássico de Chow.)

(ii) Utilizando os esquemas de Hilbert, vê-se que o funtor que, a todo $S'$
sobre $S$, associa o conjunto (reduzido a 1 ou 0 elementos) dos subesquemas
abelianos canônicos de $(\operatorname{Pic}_{X/S})\times_S S'$ é
representável por um esquema $T$ de tipo finito sobre $S$. Em virtude de (i),
$T\longrightarrow S$ é um monomorfismo e, em virtude de (3.2), é
sobrejetivo. Dizer que existe um subesquema abeliano canônico de
$\operatorname{Pic}_{X/S}$ significa que $T$ tem uma seção sobre $S$, ou ainda
que $T\longrightarrow S$ é um isomorfismo. Isso equivale também a dizer que
$T\longrightarrow S$ é étale ou, no caso em que $S$ é reduzido, que
$T\longrightarrow S$ é próprio. Daí resulta imediatamente (ii).

(iii) Utilizando simplesmente a definição de $\operatorname{Pic}_{X/S}$,
constata-se que, para todo esquema abeliano $C$ sobre $S$, o dado de um
$S$-morfismo de $X$ em um espaço homogêneo principal sob $C$ equivale ao dado
de um homomorfismo de grupos
$C'\longrightarrow\operatorname{Pic}_{X/S}$, em que $C'$ é o esquema abeliano
dual de $C$. Ora, se existe o subesquema abeliano canônico $A$ de
$\operatorname{Pic}_{X/S}$, esses homomorfismos fatoram-se necessariamente por
$A$ (como se pode ver ainda utilizando os pontos de ordem finita). Daí resulta
imediatamente (iii).

\result{OBSERVAÇÕES}{3.4} Denotaremos por
$\operatorname{Pic}_{X/S}^{\circ\circ}$ o subesquema abeliano canônico de
$\operatorname{Pic}_{X/S}$, caso exista. Infelizmente, nem sempre é esse o
caso, como se pode ver fazendo variar infinitesimalmente a superfície de Igusa
(pela deformação modular de primeira ordem). É possível, contudo, que
$\operatorname{Pic}_{X/S}^{\circ\circ}$ exista ao menos se $S$ é reduzido ou,
o que equivale ao mesmo segundo (ii), se $S$ é o espectro de um anel de
valorização discreta. Sejam então $X_0,X_1$ a fibra especial e
genérica\footnote{O francês mantido traz no singular «la fibre spéciale et
générique». Os testemunhos inglês “the special and generic fibers” e espanhol
«las fibras especial y genérica» exigem o plural; o ramo corrente usa «as
fibras» e o diplomático preserva «a fibra».} de $X$, e seja
\[
A_1=\operatorname{Pic}_{X_1/K}^\circ,
\]
em que $K$ é o corpo de frações do anel de valorização $V$. Em virtude de
Koizumi, existe um esquema abeliano $A$ sobre $S$, essencialmente único, cuja
fibra genérica é $A_1$, e vê-se facilmente, como em (2.1), (i) (supondo
doravante $X$ liso sobre $S$), que o morfismo identidade de $A_1$ se prolonga em um
morfismo
\[
A\longrightarrow\operatorname{Pic}_{X/S}.
\]
Daí um homomorfismo
\[
\tag{*}A_0\longrightarrow\operatorname{Pic}_{X_0/k}^{\circ\circ},
\]
que se demonstra facilmente ser um homomorfismo sobrejetivo cujo núcleo é um
$p$-grupo finito, em que $p$ é o expoente característico do corpo residual $k$
\src{239}
(ainda utilizando os pontos de ordem finita). Isto posto, as condições
seguintes sobre $X/S$ são equivalentes:

\medskip
\noindent a. O homomorfismo precedente (*) é um isomorfismo (o que Shimura
exprimiria dizendo que a formação da «variedade de Picard» é «compatível com a
especialização»).

\noindent b. $\operatorname{Pic}_{X/S}^{\circ\circ}$ existe (então não será
outro senão $A$).

\noindent c. (Para lembrar) Os
$\operatorname{Pic}_{X_n/S}^{\circ\circ}$ existem.

Segundo a observação feita sobre o núcleo de (*), a condição (a) é satisfeita
se a característica residual é nula, mas esse resultado será generalizado de
modo notável em (3.5).

Naturalmente, se $\operatorname{Pic}_{X/S}$ é liso sobre $S$ nos pontos de
$\operatorname{Pic}_{X/S}^\circ$, este último é
aberto em $\operatorname{Pic}_{X/S}$ (cf. (1.7)) e, portanto, munido da
estrutura induzida, é um subesquema abeliano de
$\operatorname{Pic}_{X/S}$, logo igual a
$\operatorname{Pic}_{X/S}^{\circ\circ}$, que existe nesse caso. Mas há muito
mais:

\result{TEOREMA}{3.5} Sob as condições de (3.3), seja $s\in S$ tal que
$\operatorname{Pic}_{X_s/k(s)}$ seja liso sobre $k(s)$ (ou, o que equivale ao
mesmo,
\[
\dim\operatorname{Pic}_{X_s/k(s)}=\dim H^1(X_s,\mathcal O_{X_s})).
\]
Então existe uma vizinhança aberta $U$ de $s$ tal que
$\operatorname{Pic}_{X/S}$ seja liso sobre $S$ em os pontos de
$\operatorname{Pic}_{X/S}^\circ|U$\footnote{O francês mantido traz a
construção impossível «en les points». O ramo corrente verte a leitura
corrigida «aux points» por «nos pontos»; o diplomático preserva literalmente
«em os pontos».}, que é, portanto, um subesquema abeliano aberto em
$\operatorname{Pic}_{X/S}|U$. A fortiori,
$\operatorname{Pic}_{X|U/U}^{\circ\circ}$ existe.

Indiquemos o princípio da demonstração. O que precede reduz-nos ao caso em que
$S$ é o espectro de um anel artiniano local $A$, e raciocina-se por indução
sobre a ordem infinitesimal de $A$. Pode-se, portanto, supor que
$\operatorname{Pic}_{X_n/A_n}^\circ$ seja liso sobre $A_n$ e limitar-se a
provar que $\operatorname{Pic}_{X_{n+1}/A_{n+1}}^\circ$ é liso sobre
$A_{n+1}$. Constata-se que, para isso, basta construir um esquema abeliano
$P_{n+1}$ sobre $A_{n+1}$ que prolongue
\[
P_n=\operatorname{Pic}_{X_n/A_n}^\circ,
\]
e, ao mesmo tempo, um Módulo invertível $\mathcal L_{n+1}$ sobre
$X_{n+1}\times_{A_{n+1}}P_{n+1}$ que prolongue o Módulo invertível
$\mathcal L_n$ sobre $X_n\times_{A_n}P_n$ que intervém na definição do esquema
de Picard $\operatorname{Pic}_{X_n/A_n}$ como solução de um problema universal.
(N.B. --- Podemos supor que $X$ seja munido de uma seção sobre $S$~...) Para
realizar a construção, deve-se utilizar o seguinte resultado-chave: todo
esquema abeliano definido sobre um quociente de anel artiniano local pode ser
prolongado (em outros termos, o «esquema formal dos Módulos» absoluto ([3], II)
para um esquema abeliano sobre um corpo algebricamente fechado é liso sobre o
anel dos
\src{240}
vetores de Witt sobre $k$); esse próprio resultado é obtido utilizando as
propriedades formais gerais da obstrução ao levantamento e as operações do
grupo. Admitido esse resultado, começa-se por prolongar $P_n$ de uma maneira
qualquer a $P_{n+1}$; encontra-se então uma obstrução ao levantamento de
$\mathcal L_n$, obstrução que pertence a
\[
H^2(X_0\times P_0,\mathcal O_{X_0\times P_0})\otimes_k V
\qquad(\text{onde }V=\mathfrak m^{n+1}/\mathfrak m^{n+2}),
\]
e até mesmo ao subespaço
\[
H^1(X_0,\mathcal O_{X_0})\otimes
H^1(P_0,\mathcal O_{P_0})\otimes_k V
\]
(tendo em conta que a restrição de $\mathcal L_n$ aos dois fatores $X_n$ e
$P_n$ é trivial). Ora, este último espaço não é outro senão
\[
H^1(P_0,\mathcal T_{P_0/k})\otimes V,
\]
em que $\mathcal T_{P_0/k}$ é o feixe tangente a $P_0/k$; portanto, é também o
espaço que exprime a indeterminação que havia no levantamento de $P_n$ a
$P_{n+1}$ ([3], II). Por conseguinte, pode-se corrigir esse levantamento (de
uma e uma só maneira, como convém) de modo a anular a obstrução ao levantamento
de $\mathcal L_n$.

\result{COROLÁRIO}{3.6} Sob as condições de (3.5),
$R^1f_*(\mathcal O_X)$ é um Módulo localmente livre sobre $S$ em uma
vizinhança de $s$, e sua formação comuta com a mudança de base.

Com efeito, esse Módulo não é outro senão o Módulo tangente a
$\operatorname{Pic}_{X/S}$ ao longo da seção unidade.

\result{OBSERVAÇÃO}{3.7} Utilizando o mesmo argumento que em (3.5), pode-se
mostrar que, se $S'$ é um subesquema de $S$ definido por um Ideal coerente
nilpotente e se se supõe apenas que $\operatorname{Pic}_{X'/S'}^\circ$ existe e
tem fibras lisas, então $\operatorname{Pic}_{X/S}^\circ$ existe
necessariamente e é um esquema abeliano sobre $S$. Isso permite construir
esquemas de Picard em certos casos apesar da ausência de uma hipótese
projetiva; por exemplo, o esquema abeliano dual de um esquema abeliano sobre um
anel artiniano sempre existe. Por outro lado, utilizando (3.6) no caso em que
$X$ é um esquema abeliano sobre $S$ e utilizando a estrutura conhecida de
$H^*(X_s,\mathcal O_{X_s})$ como álgebra exterior sobre
$H^1(X_s,\mathcal O_{X_s})$ (Rosenlicht--Serre), conclui-se que, para todo
$i$, $R^if_*(\mathcal O_X)$ é localmente livre e, mais precisamente, isomorfo
à $i$-ésima potência exterior de $R^1f_*(\mathcal O_X)$.

\result{OBSERVAÇÃO}{3.8} No caso de um morfismo projetivo liso
$f:X\longrightarrow S$, com $S$ reduzido e de características residuais
nulas, o resultado (3.6) era conhecido por métodos transcendentes, como
consequência da teoria de Hodge. De fato, todos os
$R^pf_*(\Omega^q_{X/S})$ são então localmente livres. Há, em contrapartida,
contraexemplos
\src{241}
no caso de características mistas para «$R^1f_*(\mathcal O_X)$ localmente
livre», com o auxílio das variedades de Serre ([8], p.~50). Parece não haver
nenhum contraexemplo em característica igual.

\result{COROLÁRIO}{3.8} Sob as condições de (3.5),
$\operatorname{Pic}_{X/S}|U$ é liso sobre $U$ em todos os pontos de
$\operatorname{Pic}_{X/S}^\sigma|U$.

Aplica-se (2.10), (i). Em particular, tendo em conta (2.10), (iii):

\result{COROLÁRIO}{3.9} Suponhamos que $S$ tenha características residuais
nulas. Então $\operatorname{Pic}_{X/S}^\tau$ é liso sobre $S$.

\result{OBSERVAÇÃO}{3.10} Deduz-se daí, por exemplo, se
$\operatorname{Pic}_{X/S}$ é também próprio sobre $S$, que o grupo de torção
dos grupos de Néron--Severi das fibras geométricas de $f$ permanece constante
sobre toda componente conexa de $S$ (o que, aliás, é evidente por métodos
transcendentes quando $f$ é liso e projetivo). Assinalemos que a utilização
direta de (2.5) permite mostrar, mais geralmente, que, se
$\operatorname{Pic}_{X/S}$ é próprio sobre $S$ (por exemplo,
$f:X\longrightarrow S$ liso e projetivo) e se $q$ é um número primo distinto
das características residuais de $S$, então a componente $q$-primária dos
grupos de torção de Néron--Severi das fibras geométricas de $X/S$ permanece
constante sobre toda componente conexa de $S$. Contudo, já não ocorre o mesmo
em característica $p>0$ para a componente $p$-primária da torção. Continua
possível, porém, que o posto sobre $k(s)$ de
\[
\operatorname{Pic}_{X_s/k(s)}^\tau/
\operatorname{Pic}_{X_s/k(s)}^{\circ\circ}=T_{X_s/k(s)}
\]
permaneça localmente constante; quando $S$ é reduzido, pode-se mostrar que isso
equivale a que $\operatorname{Pic}_{X/S}^{\circ\circ}$ exista e que
$\operatorname{Pic}_{X/S}^\tau$ seja plano sobre $S$, e basta testar a questão
para $S$ espectro de um anel de valorização discreta. Foi isso que verifiquei
nos poucos exemplos que examinei; mas, como o enunciado correspondente com $S$
artiniano é falso (cf. observações (2.9) e (3.4)), sem dúvida não convém
alimentar ilusões excessivas.

\fsection{4}{O teorema de finitude para o esquema de Picard}

Seja $f:X\longrightarrow S$ um morfismo projetivo e plano tal que
$\operatorname{Pic}_{X/S}$ exista; então a consideração dos «polinômios de
Hilbert» $Q\in\mathbf Q[t]$ permite decompor $\operatorname{Pic}_{X/S}$ em uma
soma de abertos $\operatorname{Pic}_{X/S}^Q$. Quando não se fazem hipóteses
adicionais, que assegurem, por exemplo, que $\operatorname{Pic}_{X/S}$ seja
separado sobre $S$, não será verdade, em geral, que esses abertos sejam de tipo
finito sobre $S$; obtém-se um contraexemplo quando $X$ é «uma cônica que
degenera em duas retas». É
\src{242}
possível, contudo, que assim seja se $f$ é separável e tem fibras
geometricamente irredutíveis. A questão está ligada à de saber se
$\operatorname{Pic}_{X/S}^\tau$ é de tipo finito sobre $S$, o que talvez seja
verdade sem hipótese alguma sobre as fibras de $X/S$. Quando
$f:X\longrightarrow S$ é liso, note-se que
$\operatorname{Pic}_{X/S}^\tau$ está contido em um dos
$\operatorname{Pic}_{X/S}^Q$ (graças ao fato de que, sobre uma variedade
projetiva não singular, a equivalência «por torção» é mais fina que (na
realidade, «idêntica à», segundo Matsusaka) a equivalência numérica para os
divisores), portanto é de tipo finito sobre $S$ se os
$\operatorname{Pic}_{X/S}^Q$ o são. Note-se que as questões de finitude que
acabamos de levantar conservam um sentido evidente mesmo sem supor a existência
de $\operatorname{Pic}_{X/S}$, pois se exprimem dizendo que certas famílias de
Módulos invertíveis são «limitadas» no sentido de [3], IV: consideremos, para
todo corpo algebricamente fechado $k$, os subesquemas íntegros de
$\mathbf P_k^r$, de dimensão $n$ e de grau $d$, e os Módulos invertíveis sobre
tais pré-esquemas, que tenham polinômio de Hilbert $Q$ (onde $r,n,d,Q$ são
dados); trata-se de mostrar que a família desses Módulos (considerados como
Módulos coerentes sobre os $\mathbf P_k^r$) é limitada, isto é, pode ser
parametrizada por um esquema de tipo finito sobre $\mathbf Z$.

Utilizando o método de Matsusaka [6], uma demonstração bastante técnica (que
recorre aos «critérios de equivalência» sob uma forma conveniente) permite
responder afirmativamente quando nos restringimos às subvariedades não
singulares de $\mathbf P_k^r$. Mais precisamente, obtém-se o resultado seguinte:

\result{TEOREMA}{4.1} Sejam: $X\longrightarrow S$\footnote{A leitura impressa
é preservada neste ramo; o ramo corrente escreve «Seja
$f:X\longrightarrow S$», pois o morfismo utilizado a seguir é $f$.} um
morfismo liso e projetivo com fibras geometricamente conexas, sendo $S$
noetheriano, $\mathcal O_X(1)$ um Módulo muito amplo sobre $X$ relativamente a
$S$, $E$ um subconjunto de $\operatorname{Pic}_{X/S}$, correspondente a uma
família $(\mathcal L_i)$ de Módulos invertíveis sobre as fibras geométricas
$\overline X_{s_i}$ de $X/S$, $D_i$ um divisor (não necessariamente efetivo)
sobre $\overline X_{s_i}$ que define $\mathcal L_i$,
\[
a_n^{(i)}t^n+\cdots+a_0^{(i)}
\]
o polinômio de Hilbert de $\mathcal L_i$, $\xi=\xi_i$ um divisor que define
$\mathcal O_X(1)$, isto é, uma seção hiperplana. As condições seguintes são
equivalentes:

\noindent a. $Q$\footnote{A leitura impressa é preservada neste ramo; o ramo
corrente escreve «$E$», pois o subconjunto introduzido no enunciado, cuja
quase-compacidade exprime que a família é limitada, é $E$.} é quase-compacto,
isto é, está contido em um aberto de tipo finito sobre $S$, isto é, a família
dos $(\mathcal L_i)$ é limitada.

\noindent b. $E$ está contido na reunião de um número finito dos conjuntos
$\operatorname{Pic}_{X/S}^Q$, $Q\in\mathbf Q[t]$, isto é, o conjunto dos
polinômios de Hilbert dos $\mathcal L_i$ é finito.

\noindent b'. (Se as fibras de $X/S$ têm todas a mesma dimensão $n$.) Os dois
coeficientes $a_{n-1}^{(i)}$ e $a_{n-2}^{(i)}$ dos polinômios de Hilbert dos
$\mathcal L_i$ permanecem em um conjunto finito.
\src{243}
\noindent b''. (Se as fibras de $X/S$ têm todas a mesma dimensão $n$.) Os
inteiros
\[
\xi^{n-1}D_i\qquad\text{e}\qquad\xi^{n-2}D_i^2
\]
(graus projetivos de $D_i$ e de $D_i^2$) permanecem majorados.

\result{COROLÁRIO}{4.2} Seja $f:X\longrightarrow S$ um morfismo liso e
projetivo com fibras geometricamente conexas; então os esquemas
$\operatorname{Pic}_{X/S}^Q$ e $\operatorname{Pic}_{X/S}^\tau$
($Q\in\mathbf Q[t]$) são projetivos sobre $S$.

Como são de tipo finito sobre $S$ segundo (4.1), pode-se aplicar (2.4).

\section*{BIBLIOGRAFIA}
\addcontentsline{toc}{section}{Bibliografia}

\begin{description}
\item[{[1]}] CHEVALLEY (Claude). -- Sur la théorie de Picard,
\emph{Amer. J. of Math.}, t. 82, 1960, p. 435--490.

\item[{[2]}] GROTHENDIECK (A.) e DIEUDONNÉ (J.). -- Éléments de géométrie
algébrique, Capítulo I e seguintes. -- Paris, Presses universitaires de France,
1960--1961 (Institut des Hautes Études Scientifiques, 4, 8, 11, \ldots).

\item[{[3]}] GROTHENDIECK (Alexander). -- Technique de descente et théorèmes
d'existence en géométrie algébrique, I--V, \emph{Séminaire Bourbaki}, t. 12,
1959/60, n\textsuperscript{o} 190, 29 p. e n\textsuperscript{o} 195, 22 p.;
t. 13, 1960/61, n\textsuperscript{o} 212, 20 p. e
n\textsuperscript{o} 221, 28 p.; t. 14, 1961/62,
n\textsuperscript{o} 232, 19 p.

\item[{[4]}] GROTHENDIECK (Alexander). -- Séminaire de géométrie algébrique,
1960/61. -- Paris, Institut des Hautes Études Scientifiques, 1961.

\item[{[5]}] GROTHENDIECK (Alexander). -- Séminaire de géométrie algébrique,
1961/62, redigido por um grupo de ouvintes. -- Paris, Institut des Hautes Études
Scientifiques (a publicar).

\item[{[6]}] MATSUSAKA (T.). -- The criteria for algebraic equivalence and the
torsion group, \emph{Amer. J. Math.}, t. 79, 1957, p. 53--66.

\item[{[7]}] MUMFORD (D.) e TATE (J.). -- Séminaire de géométrie algébrique,
Harvard University, semestre da primavera de 1962 (a publicar).

\item[{[8]}] SERRE (Jean-Pierre). -- Sur la topologie des variétés algébriques
en caractéristique $p$, \emph{Symposium internacional de topología algebraica}
[1956, México], p. 24--53. -- México, Universidad Nacional Autónoma, 1958.
\end{description}
